\section{Complete supports and original-action information}
\label{supp:fullsupport125}

\begin{proof}[Proof of Theorem~\ref{thm:fullinterpolation125}]
For necessity let $q$ and its supporting slopes be compatible. The
function $h=q-\mu x^2/2$ is differentiable with a nondecreasing
$K$-Lipschitz derivative. On an adjacent interval the derivative
increment $r(t)=h'(x_i+t)-g_i$ starts at zero, ends at $s_i$, and
satisfies
\[
 \max\{0,s_i-K(d_i-t)\}\le r(t)\le\min\{Kt,s_i\}.
\]
Its integral is $w_i$. Integrating the two ramps gives exactly
\eqref{eq:adjacentfull125}. If $K=0$, the derivative is constant,
and $s_i=w_i=0$ instead. At every feasible displacement $t$ from
$x_i$, the smooth upper bound and the global support give
\[
 \ell_j(x_i;p_j)+(p_j-\mu a_j)t
 \le h(x_i+t)\le f_i+g_it+Kt^2/2.
\]
Rearranging and maximizing over $[-x_i,A-x_i]$ proves
\eqref{eq:obstaclefull125}. This step uses the same $p_j$ at
every node and displacement. All paid intervals and the integrated
derivative-cap boxes hold by construction.

For sufficiency first take $K>0$. Fix an adjacent interval and
temporarily suppress its index. Thus $d>0$, $0\le s\le Kd$, and
$s^2/(2K)\le w\le sd-s^2/(2K)$. If $D=d-s/K>0$, define
\[
 c=\frac{w-s^2/(2K)}{D}\in[0,s],\qquad
 r(t)=\min\{Kt,c\}+\max\{0,s-c-K(d-t)\}.
\]
The first ramp ends at $c/K$, and the last starts at
$d-(s-c)/K$. Their ordered endpoints follow from $s\le Kd$.
The function $r$ is continuous, nondecreasing and $K$-Lipschitz;
its endpoint values are $0,s$. Direct integration yields
\[
 \int_0^d r(t)\,dt
   =cd-\frac{c^2}{2K}+\frac{(s-c)^2}{2K}
   =cD+\frac{s^2}{2K}=w.
\]
Set
$h(x_i+t)=f_i+g_it+\int_0^t r(u)\,du$.
If $D=0$, then $s=Kd$ and both allowed area endpoints equal
$Kd^2/2$; set $r(t)=Kt$. These definitions realize both endpoint
values and derivatives. Joining the pieces gives a continuously
differentiable function on $[0,A]$. Its derivative increments are
between zero and $K$ times distance within each piece and, by
addition, across joins.

It remains to verify every affine support everywhere. On one
adjacent interval write an affine support as
$\ell(x_i+t)=\ell(x_i)+pt$, and let
$G_0=f_i-\ell(x_i)$ and
$G_1=f_{i+1}-\ell(x_{i+1})$ be the endpoint gaps.
If $p\le g_i$, the gap has nonnegative derivative throughout and
its minimum is $G_0\ge0$. If $p\ge g_{i+1}$ its minimum is
$G_1\ge0$. In the remaining case put $v=p-g_i\in(0,s)$.
For $D>0$, if $v\le c$, the derivative of the gap first reaches
zero at $t=v/K$ on the first ramp; its minimum is
\[
                         G_0-v^2/(2K).
\]
If $v\ge c$, it reaches zero at
$t=d-(s-v)/K$ on the last ramp, and its minimum is
\[
                         G_1-(s-v)^2/(2K).
\]
At $v=c$ the gap is constant between the ramps and the two
expressions agree, by the area identity above. Because
$v/K\le d$ and $(s-v)/K\le d$, both displacements are feasible
in the full domain from the corresponding endpoint. The two node
inequalities \eqref{eq:obstaclefull125} therefore make both
displayed minima nonnegative. For $D=0$ the derivative is the
single $K$-ramp, and the first expression proves the same conclusion.
The endpoint-slope cases $v=0,s$ are already covered. Thus the
constructed function lies above every support on every interval,
using the fixed shared slopes in the proposed list. No assertion
about an unconstrained optimizer is used in this construction.

For $K=0$, the adjacent equalities join to one affine function
$h=f_0+g_0x$. For a support of slope $p$, its global minimum gap
equals the value at the appropriate domain endpoint. Equivalently
the exact condition at any node is
$f_i-\ell(x_i)\ge\max\{(p-g_i)(-x_i),(p-g_i)(A-x_i)\}$,
which is \eqref{eq:obstaclefull125}. Therefore this case is also
necessary and sufficient without division by $K$.

Adding back $\mu x^2/2$ gives the asserted strongly convex,
$L$-smooth $q$. Its derivative is nondecreasing and lies between
the original $q'$ values at the endpoints of each piece. Thus the
node derivative cap implies the cap everywhere. The paid intervals
are realized at their nodes, and every original quadratic support
holds globally. This proves sufficiency.

Each adjacent inequality defines a convex set: its nonaffine term
is the square of an affine function. For fixed feasible displacement
$t$, the expression
$(p_j-\mu a_j-g_i)t-Kt^2/2$
is affine in the optimization variables. Its supremum is convex,
and the left side of \eqref{eq:obstaclefull125} is affine.
Consequently all support inequalities are convex even when the
support slopes are variables. All constraints are closed. The
finite value, derivative and supplied slope boxes make the feasible
set compact; compatibility makes it nonempty. Linear contrasts
therefore have attained extrema. Necessity and constructive
sufficiency identify these extrema with the function-class extrema,
including the deterministic curvature adjustment when returning from
$h$ to $q$. Removing constraints enlarges this class and proves the
claimed dominance.
\end{proof}

\begin{proof}[The residual certificate \eqref{eq:fullresidual125}]
For every feasible $z$, $H(z)\ge H(z)+\sum_jy_jc_j(z)$
because $y_j\ge0$ and $c_j(z)\le0$. Convexity and the chosen
subgradients give
\[
 H(z)+\sum_jy_jc_j(z)\ge \mathcal L_0+R^{\mathsf T}(z-z_0).
\]
The minimum of this affine expression over the enclosing box is
exactly $B$. Hence $H(z)\ge B$ for every feasible point and
$q(b)-q(a)\le-B$ for every compatible original objective.
For $K>0$ the maximizing displacement in
\eqref{eq:boundedconjugate125} is unique, so differentiation gives
the clipped expression in the main text; its values agree at the
two clipping transitions. For $K=0$ choose the active endpoint
as a subgradient of $\max\{v\ell,vr\}$, and either endpoint or
any convex combination when $v=0$. The chain rule for affine
arguments supplies a valid subgradient of each support constraint.
All these choices can be rational. Constraint values, subgradients,
the box residual and the final bound are then exact rational
operations. No feasibility or stationarity assumption on $z_0$
has entered the proof.
\end{proof}

For completeness, $A=0$ is an identity problem: the only feasible
action is zero and every action contrast is exactly zero. One
checks the paid value intervals and the supports at that point;
neither derivative interpolation nor an action query is required.
Repeated anchors do not create zero-length intervals. An empty
feasible set indicates inconsistent supplied data and must be
reported; constraints cannot be discarded to manufacture a gain.

\subsection{Interval slopes and a global robust lower obstacle}

The shared slope variables are an exact convex lifting, not additional
information relative to enforcing the entire robust pointwise obstacle.
To see this, let a convex function $h$ satisfy $h(a)\ge W$, and
suppose
\[
 h(x)\ge W+\min_{p\in[p^-,p^+]}p(x-a)
                 \qquad(0\le x\le A).
\]
Define
\[
 L_a=\sup_{x<a}\frac{h(x)-W}{x-a},\qquad
 R_a=\inf_{x>a}\frac{h(x)-W}{x-a},
\]
with the empty-side conventions $L_a=-\infty$ and $R_a=+\infty$.
For $x<a<y$, convexity and $W\le h(a)$ give
$(h(x)-W)/(x-a)\le(h(y)-W)/(y-a)$, hence $L_a\le R_a$.
The robust obstacle gives $L_a\le p^+$ and $R_a\ge p^-$.
Consequently $[L_a,R_a]\cap[p^-,p^+]$ is nonempty. Any $p$
in that intersection satisfies $h(x)\ge W+p(x-a)$ everywhere.
The converse follows by taking the minimum over the slope interval.
This also proves the endpoint cases with the empty-side conventions.
Adding back the common quadratic establishes the assertion for
\eqref{eq:globalpayload125}.

\subsection{An exact action contrast that uses the omitted support}
\label{supp:fullsupportwitness125}

\begin{proposition}[A strict economic effect of retaining a paid support]
\label{prop:fullsupportwitness125}
Take $A=1$, $\mu=1$, $L=2$, $M=3$, and disclose
\[
 q(0)=0,\quad q(1/2)\in[0,1],\quad
 q'(1/2)\in[-5/4,3].
\]
The paid support at $a_1=1/2$ has $W_1=0$ and exact slope
$p_1=1$, so
\[
              S_1(x)=x^2/2+x/2-3/8.
\]
Let $\mathcal Q_{\rm red}$ be the class retaining the displayed
levels and derivatives and the curvature and derivative-cap assumptions,
but omitting $q\ge S_1$. Let $\mathcal Q_{\rm full}$ also retain
that support. For the incumbent $a=1$ and candidate $b=0$,
\[
 \max_{q\in\mathcal Q_{\rm red}}[q(0)-q(1)]=-1/4,
 \qquad
 \max_{q\in\mathcal Q_{\rm full}}[q(0)-q(1)]=-5/8.
\]
Both maxima are attained. With arbitrary Borel
$\eta:[0,1]\to[0,\delta]$, the corresponding maximum original-policy
contrasts are $-1/4+\delta$ and $-5/8+\delta$.
\end{proposition}

\begin{proof}
Strong convexity at the midpoint gives
$q(1/2)\le[q(0)+q(1)]/2-1/8$, hence $q(1)\ge1/4$.
The function $q_{\rm red}(x)=x^2/2-x/4$ attains this bound.
It has curvature one, derivative between $-1/4$ and $3/4$,
$q_{\rm red}(1/2)=0$, and $q_{\rm red}'(1/2)=1/4$; it therefore
satisfies every reduced-class restriction. In the full class the
support gives $q(1)\ge S_1(1)=5/8$. The function
$q_{\rm full}(x)=x^2/2+x/8$ attains this bound. Its midpoint
value is $3/16$, its midpoint derivative is $5/8$, and its
derivative remains within the cap. Moreover
\[
                  q_{\rm full}(x)-S_1(x)=3(1-x)/8\ge0.
\]
Thus it satisfies the entire global support. Both functions vanish
at zero, proving the two contrast maxima. Taking $\eta(0)=\delta$
and $\eta(1)=0$ attains the upper future-error difference for each
witness, while the uniform bound gives the reverse inequality.
\end{proof}

The derivative interval in this example does not hide a stronger free
derivative query. Converting the support with its paid upper level
$U_1=1$ uses gap $e=1$, $K=1$ and both feasible radii $1/2$.
The exact bound is $\Psi(1,1/2,1/2)=9/4$, giving
$[-5/4,13/4]$, and intersection with $[-3,3]$ gives the stated
interval. The improvement is thus over the level--derivative summary of
that payload. Directly evaluating $S_1(1)$ already obtains the stronger
answer here; no exclusive benefit of a conic solver or neural
representation is asserted. If $1/4\le\delta<5/8$, this omitted
information changes whether the proposed policy change is uniformly
certified at all.

\subsection{Why nodewise obstacle checks are insufficient}

Take $K=1$ and $h(x)=x^2/2-x/2$ on $[0,1]$, with nodes
$0,1/4,1$. Its jets satisfy every adjacent interpolation condition.
The affine obstacle $\ell(x)=-7/64$ lies below the function at
these nodes: their gaps are $7/64,1/64,7/64$. Nevertheless
$h(1/2)=-1/8<-7/64$. At the middle node $g=-1/4$, so
$v=1/4$ and
\[
 \Phi_1(1/4;-1/4,3/4)=1/32>1/64=f-\ell.
\]
The new constraint rejects this impossible global-support payload.
Equivalently, with $\mu=1$, the proposed objective is
$q(x)=x^2-x/2$ and the proposed support is
$S(x)=x^2/2-7/64$. It can be written at anchor $1/4$ with
$W=-5/64$ and slope $p=1/4$. This is a negative algebraic check
of inconsistent data, not a claimed admissible economic observation.

\section{Proof of the complete acquisition account}
\label{supp:paidresolution125}

\begin{proof}[Initial bracket and Theorem~\ref{thm:paidwork125}]
Let $a_q^*$ be the constrained minimizer. Its variational inequality
and strong convexity give
$q(a_0)-q(a_q^*)\ge\mu|a_0-a_q^*|^2/2$.
Since $q(a_0)\le J(a_0)$, the original-policy regret bound implies
\eqref{eq:warmbracket125}, including either capacity boundary.
An outward radius and intersection with another valid optimizer
interval cannot exclude $a_q^*$.

Write $w$ for the actual interval width at one planned step. If
$w\le t$, no query is necessary. Otherwise the actual midpoint's
distance from each domain boundary is at least
$(1/2-\xi)w>(1/2-\xi)t\ge h$. Thus both trial displacements
$\pm h$ in the support conversion are feasible. Smoothness and
the support inequality imply
\[
 p^- - e/h-Kh/2\le q'(a)\le p^+ +e/h+Kh/2.
\]
The interval width is at most $v+2e/h+Kh$. Under
\eqref{eq:jointprecision125} its three terms are at most
$\mu t/4,\mu t/2,\mu t/4$, respectively. For $K=0$ the
last term is zero. The derivative interval therefore has width
at most $\mu t$. The inherited rounded-midpoint result now gives
new optimizer width at most
\[
 \max\{\gamma w,\mu t/\mu\}\le
 \max\{\gamma w_k,t\}=t.
\]
This argument also covers a sign-resolved derivative interval,
which may contract at a coarser realized precision than requested.
It uses a valid original derivative enclosure; differentiating the
fitted critic does not establish the premise.

Minimality of $N$ gives $\gamma w_0\le w_1<w_0$, with the
upper inequality weak if a redundant boundary case is retained.
Subsequent planned widths satisfy $w_{j+1}=\gamma w_j$.
Induction using the preceding step proves the asserted final width.
Summing the per-probe work bounds gives
\[
 \sum_{j=1}^N(C_0+C_1w_j^{-p})
 =NC_0+C_1w_*^{-p}\sum_{j=0}^{N-1}\gamma^{pj},
\]
which is \eqref{eq:geometricwork125}. Skipped steps can only reduce
this nonnegative bound. For the certificate-planning optimality
claim, any $N$-step plan reaching $w_N\le w_*$ under
$w_{j+1}\ge\gamma w_j$ must have
$w_j\le w_*\gamma^{-(N-j)}$. The displayed schedule takes the
largest feasible width at every step, minimizing each nonincreasing
cost term simultaneously. Fewer than $N$ steps cannot attain the
target under the guaranteed contraction floor. A longer plan has
the same upper constraints on its final $N$ widths; its additional
nonnegative cost terms cannot improve the minimum. This proves
optimality for this declared bound and planning interface, without
asserting that an actual oracle must cost the upper bound or that
observations cannot produce better contraction.

Finally $a_q^*$ lies in the certified interval, giving
$|b-a_q^*|\le r_b$. Integrating the derivative cap gives
$q(b)-q(a_q^*)\le Mr_b$. Adding $0\le\eta(b)\le\delta$
proves the regret bound in \eqref{eq:terminalpolicy125} directly,
without an absolute-value width penalty. Also
$q(a_q^*)\ge q(b)-Mr_b\ge W_b-Mr_b$ supplies a lower optimum
bound when one is needed. Combining only this latter lower bound
with $U_b+\delta$ would give the weaker bound
$e_b+\delta+Mr_b$; an implementation can use the minimum of
all valid regret certificates instead. If $q'(a_q^*)=0$,
the smooth Taylor upper inequality instead gives
$q(b)-q(a_q^*)\le Lr_b^2/2$ and the corresponding replacement.
An interval strictly inside the feasible domain contains an
interior minimizer and therefore certifies that stationary condition.
At a capacity minimum its derivative can be nonzero, which is why
the quadratic bound is not asserted there without further evidence.
The absolute interval for $J$ has width $e_b+\delta$, a separate
condition from optimizer localization. The payment for its value
query and for all other service operations is not present in the
sum over the $N$ localization probes and must be added explicitly.
\end{proof}

\subsection{Sharp terminal information requirements}

The following exact reduced-interface statement separates action
resolution from the precision of an absolute policy value. It does
not claim a lower bound after raw supports, additional value data or
the original investment formulas have also been supplied.

\begin{proposition}[Matching terminal regret at a stationary bracket interface]
\label{prop:terminalminimax125}
Suppose the only action information is $0\le\ell\le r\le A$,
$0<\mu\le L$, and the fact that the unique minimizer $z$ belongs
to $[\ell,r]$ and satisfies $q'(z)=0$. Let $\mathcal Q_B$ contain
all differentiable $\mu$-strongly convex functions with
$L$-Lipschitz derivative and these properties. A derivative cap,
if imposed for this statement, is at least $LA$.
For every Borel law $\nu$ of $B\in[0,A]$,
\begin{align*}
 \sup_{q\in\mathcal Q_B}\mathbb E_\nu[q(B)-\min q]
 =\frac L2\left\{\operatorname{Var}_\nu(B)+
     \max\big[(\mathbb E_\nu B-\ell)^2,
              (\mathbb E_\nu B-r)^2\big]\right\}.
\end{align*}
Hence
\[
 \inf_\nu\sup_{\substack{q\in\mathcal Q_B\\0\le\eta\le\delta}}
 \mathbb E_\nu[q(B)+\eta(B)-\min q]
                   =\delta+L(r-\ell)^2/8.
\]
The deterministic midpoint attains the bound; randomization cannot
improve it. Independently, if the only absolute-level information at
a chosen action is $q(b)\in[W,U]$, with $0\le\eta(b)\le\delta$
and arbitrary additive shifts otherwise allowed, the exact compatible
range of $J(b)$ is $[W,U+\delta]$.
\end{proposition}

\begin{proof}
Smoothness at the stationary minimizer gives
$q(B)-q(z)\le L(B-z)^2/2$. Taking expectations and maximizing
over $z\in[\ell,r]$ bounds the left side by the stated expression,
because a convex quadratic in $z$ has its interval maximum at an
endpoint. For either endpoint $z$ the single function
$q_z(x)=L(x-z)^2/2$ belongs to the class, is also
$\mu$-strongly convex, satisfies the permitted cap, and attains
its corresponding expression. Choosing the maximizing endpoint
proves equality; the adverse function knows the law before the
random action is drawn. The expression is minimized by zero
variance and mean $(\ell+r)/2$, giving $L(r-\ell)^2/8$.
The constant error $\eta\equiv\delta$ attains the additional
worst-case allowance. For the absolute-level statement, an arbitrary
compatible shape can be shifted to give either endpoint $W,U$ at
$b$. Taking the policy error respectively zero or $\delta$ attains
both endpoints of $[W,U+\delta]$; the assumed bounds give its
containment. No optimizer localization can resolve this free level
shift without further value information.
\end{proof}

